\documentclass[10pt,a4paper]{amsart}
\usepackage[margin=19mm]{geometry}
\usepackage{amsmath,amssymb,mathtools}
\usepackage[scr=rsfs,cal=euler]{mathalfa}
\usepackage{xfrac}
\usepackage[unicode,hidelinks,bookmarks=false]{hyperref}
\newcommand{\ie}{i.e.\ }
\newcommand{\cf}{cf.\ }
\newcommand{\resp}{respectively\ }
\newcommand{\Subsection}{\subsection}
\providecommand{\nobreakdash}{\nobreak\hbox{-}}
\setlength{\emergencystretch}{2em}
\multlinegap0pt
\usepackage{bm}
\usepackage{bbm}
\usepackage{dsfont}
\usepackage{xspace}
\usepackage{cancel}
\usepackage{mathtools}
\usepackage{amsthm}
\makeatletter
\@ifundefined{Theorem}{\newtheorem{Theorem}{Theorem}}{}
\@ifundefined{Lemma}{\newtheorem{Lemma}[Theorem]{Lemma}}{}
\makeatother
\def \H {{\mathbb H}}
\def \L {\mathcal{L}}
\newcommand{\J}{\mathcal{J}}
\newcommand{\R}{{\mathbb R}}
\title[Boundary regularity for subelliptic equations in the Heisenb]{Boundary regularity for subelliptic equations in the Heisenberg group}
\author{Farhan Abedin, Giulio Tralli}
\date{Source: arXiv:2506.05151v1 (CC BY 4.0)}
\begin{document}
\maketitle

\noindent\emph{Source and licence.} Boundary regularity for subelliptic equations in the Heisenberg group. Version arXiv:2506.05151v1 (CC BY 4.0), distributed under the Creative Commons Attribution 4.0 International licence, as recorded in the arXiv licence metadata for this exact version and shown on the abstract page https://arxiv.org/abs/2506.05151v1. Full rights notice: licenses/arxiv-2506.05151-CC-BY-4.0.txt.\par\smallskip

\noindent\textit{Editorial scope.} Counted unit: the Theorem whose source label is \texttt{primadinocond} in the article (source0.tex lines 930--939, source label \texttt{primadinocond}), delivered together with the article's own complete original proof of it (source0.tex lines 940--1024, one \texttt{proof} environment of 85 lines). No mathematics is written, repaired or re-derived: statement and proof are the article's own text line for line. Required context retained uncounted: lemprimogiro (lines 881--888). Every definition, display and label of those lines is retained unchanged. Omitted from this package and not redistributed: 1--880, 889--929, 1025--1344. Nothing inside a retained excerpt is omitted or altered: every retained body is delivered exactly as the article's lines are written, so no omission receipt of any kind accompanies this package. Citations: the retained excerpts carry no citation command at all, so no bibliography entry of the article is transcribed and no reference is redistributed. Reference adaptation: none was needed and none was made; every reference inside the delivered unit resolves inside it, so no unresolved reference or citation is produced. Printed slips kept literal: no printed word, symbol, hypothesis or number of the counted statement or of its proof was altered. Layout and class: the article's own class is \texttt{amsart} with packages including amsmath, amsfonts, amssymb, amsthm, graphicx, hyperref, color, geometry; the wrapper is the lane's pinned \texttt{amsart} editorial wrapper at 10pt on A4, which restates the theorem-like environments the retained text needs and adds the article's own macro definitions that the retained text needs (\texttt{\textbackslash H}, \texttt{\textbackslash J}, \texttt{\textbackslash L}, \texttt{\textbackslash R}), with extra wrapper packages (bm, bbm, dsfont, xspace, cancel, mathtools). The delivered numbering is the unit's own. Assets: no retained span contains a figure environment, a graphics-inclusion command, a PDF-page inclusion, a table or a TikZ picture; no image file is copied into this package, the package declares no asset dependency and no image file is redistributed with it. Licence: version arXiv:2506.05151v1 of this article is distributed under the Creative Commons Attribution 4.0 International licence, as recorded in the arXiv licence metadata for this exact version and shown on the abstract page https://arxiv.org/abs/2506.05151v1. A full rights notice is delivered at licenses/arxiv-2506.05151-CC-BY-4.0.txt. Characters: every retained excerpt is pure ASCII, so the delivered engine, XeLaTeX with its default Latin Modern text font, reports no missing character.
\begin{Lemma}\label{lemprimogiro}
Fix $r_0>0$. Suppose $u\in C^2(B_{4r_0}(0)\cap\H^n_+)\cap C(\overline{B_{4r_0}(0)\cap\H^n_+})$ satisfies $\L_A u \geq f$ in $B_{4r_0}(0)\cap\H^n_+$ for $f \in L^{\infty}(B_{4r_0}(0)\cap\H^n_+,|x|^2)$, and $u \leq 0$ on $B_{4r_0}(0) \cap \{t=0\}$. Then
$$u(z) \leq \frac{M_0}{r_0^4} \left(|x|^4+r_0^2 t \right) \quad \text{ for all } \ z \in B_{\frac{4}{3}r_0}(0)\cap \H^n_+$$
with
$$M_0=C_0\max\left\{ ||u^+||_{L^{\infty}(B_{4r_0}(0)\cap\H^n_+)}, 
\frac{r_0^4}{\lambda}||f^-||_{L^{\infty}(B_{4r_0}(0)\cap\H^n_+,|x|^2)}\right\}$$
and $C_0$ a structural constant depending only on $Q, \frac{\Lambda}{\lambda}$.
\end{Lemma}

\begin{Theorem}\label{primadinocond}
Fix $r_0>0$. Suppose $u\in C^2(B_{4r_0}(0)\cap\H^n_+)\cap C(\overline{B_{4r_0}(0)\cap\H^n_+})$ satisfies $\L_A u \geq f$ in $B_{4r_0}(0)\cap\H^n_+$ for $f \in L^{\infty}(B_{4r_0}(0)\cap\H^n_+,|x|^2)$, and $u\leq 0$ on $B_{4r_0}(0) \cap \{t=0\}$. Then we have
\begin{equation}\label{linear-growth-estimate}
u(z)\leq \frac{M_1}{r^2_0} t \qquad \text{ for all } \ z=(x,t) \in B_{r_0}(0)\cap \H^n_+
\end{equation}
with
$$M_1=C_1\max\left\{ ||u^+||_{L^{\infty}(B_{4r_0}(0)\cap\H^n_+)}, 
\frac{r_0^4}{\lambda}||f^-||_{L^{\infty}(B_{4r_0}(0)\cap\H^n_+,|x|^2)}\right\}$$
and $C_1>0$ a structural constant.
\end{Theorem}

\begin{proof}
We start by noticing a direct consequence of Lemma \ref{lemprimogiro}:
\begin{equation}\label{atzero}
u(0,t)\leq \frac{M_0}{r^2_0} t \quad \text{ for all } 0<t<r_0^2,
\end{equation}
where $M_0$ is the positive constant in Lemma \ref{lemprimogiro}. This proves \eqref{linear-growth-estimate} at $x = 0$, so we need to prove the estimate at each $x_0\in\R^{2n}$ with $0<|x_0|<r_0$ . Again by Lemma \ref{lemprimogiro} we have
\begin{equation}\label{atx0sopra}
u(x_0,t)\leq \frac{M_0}{r^4_0} \left(t^2+r_0^2t \right )\leq \frac{2 M_0}{r^2_0} t \qquad \mbox{ for any } t \mbox{ such that } |x_0|^2\leq t<\sqrt{r_0^4 - |x_0|^4}.
\end{equation}
The previous upper bound (which is effective if $2|x_0|^4<r_0^4$) does not say anything about the range $t\in (0,\min\left\{|x_0|^2, \sqrt{r_0^4 - |x_0|^4}\right\})$. With this in mind, let us now denote
$$
D_{x_0}=\left\{(x,t)\in\H^n_+ \,:\, |x-x_0|< \frac{1}{10}|x_0|,\, 0<t< |x|^2\right\}.$$
We will apply a barrier argument to prove the bound \eqref{linear-growth-estimate} in $D_{x_0}$.

First, notice that, for any $(x,t)\in \overline{D_{x_0}}$ we have 
\begin{equation}\label{comparax}
\frac{9}{10}|x_0|\leq |x_0|-|x-x_0|\leq |x| \leq |x|+|x-x_0|\leq \frac{11}{10}|x_0|.
\end{equation}
In particular, using that $|x|<\frac{11}{10}r_0$, we recognize that $\overline{D_{x_0}}\subseteq B_{\frac{4}{3}r_0}(0)$. Hence, we can use Lemma \ref{lemprimogiro} once more to infer
\begin{equation}\label{ustasotto}
u(x,t)\leq \frac{M_0}{r_0^4}|x|^2(|x|^2+r_0^2)\leq \frac{\tilde{M}_0}{r_0^2}|x_0|^2 \qquad\mbox{ for any }(x,t)\in \overline{D_{x_0}}
\end{equation}
where $\tilde{M}_0=\frac{11^2(11^2+10^2)}{10^4} M_0$. Letting
$$
\alpha=\frac{25n}{8}\frac{\Lambda}{\lambda}
\quad\mbox{ and }\quad M=\max\left\{\frac{\tilde{M}_0}{1- e^{\frac{-\alpha}{100}}}, \frac{e^{\frac{61\alpha}{50}}}{4n\alpha\Lambda} r^2_0 ||f^-||_{L^\infty(D_{x_0})} \right\},%\frac{2\lambda}{25n^2\Lambda^2}{}e^{\frac{61n\Lambda}{16\lambda}}
$$
we can consider the function
$$
w(z)=M \frac{|x_0|^2}{r^2_0} \left( 1- e^{-\alpha\frac{t+|x-x_0|^2}{|x_0|^2}}\right).
$$
We remark that, by \eqref{comparax}, we have
$$
\begin{cases}
w(x,t) \geq M \frac{|x_0|^2}{r^2_0} \left( 1- e^{\frac{-\alpha}{100}}\right) \qquad\,\,\,\mbox{ for } |x-x_0|=\frac{1}{10}|x_0|\mbox{ and }t\geq 0  \\
w(x,t)\geq M \frac{|x_0|^2}{r^2_0} \left( 1- e^{\frac{-81\alpha}{100}}\right) \qquad\mbox{ for } |x-x_0|\leq \frac{1}{10}|x_0|\mbox{ and }t=|x|^2.
\end{cases}
$$
Therefore,
$$
w(z)\geq \tilde{M}_0 \frac{|x_0|^2}{r^2_0}\qquad\mbox{ for any }z\in \partial D_{x_0}\cap \H^n_+.
$$
If we keep in mind \eqref{ustasotto} and the fact that $u\leq 0\leq w$ on $\{t=0\}$, the previous lower bound implies
\begin{equation}\label{ustasottoalbordo}
u\leq w\quad\mbox{ on }\partial D_{x_0}
\end{equation}
On the other hand, for any $z\in D_{x_0}$, we have
\begin{align*}
\L_A w (z)&= 
\frac{2\alpha M}{r_0^2}e^{-\alpha\frac{t+|x-x_0|^2}{|x_0|^2}}\left( \text{tr}(A(z)) -\frac{2\alpha}{|x_0|^2}\langle A(z)\left(\J x + (x-x_0)\right),\left(\J x + (x-x_0)\right) \rangle\right) \\
&\leq \frac{4\alpha M}{r_0^2}e^{-\alpha\frac{t+|x-x_0|^2}{|x_0|^2}}\left( n\Lambda -\frac{\alpha\lambda}{|x_0|^2}\left|\J x + (x-x_0)\right|^2 \right) \\
& \leq \frac{4\alpha M}{r_0^2}e^{-\alpha\frac{t+|x-x_0|^2}{|x_0|^2}}\left( n\Lambda -\frac{\alpha\lambda}{|x_0|^2}\left( |\J x| - |x-x_0|\right)^2 \right)\\
&\leq \frac{4\alpha M}{r_0^2}e^{-\alpha\frac{t+|x-x_0|^2}{|x_0|^2}}\left( n\Lambda -\frac{\alpha\lambda}{|x_0|^2}\left( \frac{9}{10}|x_0| -  \frac{1}{10}|x_0|\right)^2 \right)\\
& \leq \frac{-4n\alpha\Lambda M}{r_0^2}e^{-\alpha\frac{|x|^2+\frac{1}{100}|x_0|^2}{|x_0|^2}} \\
& \leq \frac{-4n\alpha\Lambda M}{r_0^2}e^{-\alpha\frac{61}{50}},
\end{align*}
where we used that $A(z) \in M_n(\lambda, \Lambda)$, together with the definition of $\alpha$ and \eqref{comparax}. Recalling the definition of $M$, we thus have
$$
\L_A w (z)\leq -||f^-||_{L^\infty(D_{x_0})} \leq -f^-(z)\leq f(z)\leq \L_A u (z)
 \quad \text{for all } z \in D_{x_0}.$$
Thanks also to \eqref{ustasottoalbordo}, we are then able to apply the comparison principle which yields
$$
u\leq w \quad \mbox{ in }D_{x_0}.
$$
In particular,
$$
u(x_0,t)\leq M \frac{|x_0|^2}{r_0^2}\left(1-e^{-\alpha\frac{t}{|x_0|^2}}\right)\qquad\mbox{ for any }t\mbox{ such that }0<t<|x_0|^2.
$$
The convexity of the function $\sigma\mapsto e^{-\alpha\sigma}$ (which yields $1-e^{-\alpha\sigma}\leq \alpha\sigma$ for any $\sigma\in\R$) implies
\begin{equation}\label{miracle}
u(x_0,t)\leq \frac{\alpha M}{r_0^2} t\qquad\mbox{ for any }t\mbox{ such that }0<t<|x_0|^2.
\end{equation}
The combination of \eqref{atzero}-\eqref{atx0sopra}-\eqref{miracle} yields
$$
u(z)\leq \frac{M_1}{r^2_0} t \qquad \text{ for all } \ z \in B_{r_0}(0)\cap \H^n_+
$$
if we choose
$$M_1\geq\max\left\{ 2M_0, \alpha M \right\}.$$
From a brief check of the dependence of the explicit constants and from the inequality $ ||f^-||_{L^\infty(D_{x_0})}\leq ||f^-||_{L^\infty(B_{4/3\, r_0}(0))}\leq 2 r^2_0 ||f^-||_{L^{\infty}(B_{4r_0}(0)\cap\H^n_+,|x|^2)}$, we readily realize that we can take
$$M_1=C_1\max\left\{ ||u^+||_{L^{\infty}(B_{4r_0}(0)\cap\H^n_+)}, 
\frac{r_0^4}{\lambda}||f^-||_{L^{\infty}(B_{4r_0}(0)\cap\H^n_+,|x|^2)}\right\}$$
with the choice 
$$C_1=\max\left\{ 2C_0, \frac{11^2(11^2+10^2)}{10^4} \frac{\alpha C_0}{1- e^{\frac{-\alpha}{100}}}, \frac{e^{\frac{61\alpha}{50}}}{(Q-2)\frac{\Lambda}{\lambda}} \right\}$$
where $C_0$ is the structural constant coming from Lemma \ref{lemprimogiro} and $\alpha=\frac{25}{16}(Q-2)\frac{\Lambda}{\lambda}$ as fixed above.
\end{proof}

\end{document}
